\documentclass[../main]{subfiles}
\begin{document}
\ifSubfilesClassLoaded{\setcounter{page}{7}}{}
\section{Предисловие}
\label{sec:02_predislovie}

Критикуя Готскую программу, Маркс писал: «Мы имеем здесь дело не с таким коммунистическим обществом, которое развилось на своей собственной основе, а, напротив, с таким, которое только что выходит как раз из капиталистического общества и которое поэтому во всех отношениях, в экономическом, нравственном и умственном, сохраняет еще родимые пятна старого общества, из недр которого оно вышло» и «На высшей фазе коммунистического общества, после того как исчезнет порабощающее человека подчинение его разделению труда; когда исчезнет вместе с этим противоположность умственного и физического труда; когда труд перестанет быть только средством для жизни, а станет сам первой потребностью жизни; когда вместе с всесторонним развитием индивидов вырастут и производительные силы и все источники общественного богатства польются полным потоком, лишь тогда можно будет совершенно преодолеть узкий горизонт буржуазного права, и общество сможет написать на своем знамени: Каждый по способностям, каждому по потребностям!».

Идее частной собственности и господства капитала противостоит идея коммунизма. Обществу, где господствует капитал, в реальности можно противопоставить только коммунистическое общество. Но что такое коммунизм?

В этой брошюре нас интересует то, что Марксом было названо высшей фазой коммунизма, то есть развитое на своей собственной основе коммунистическое общество, или коммунизм как таковой. Низшая фаза коммунизма, или переходный период от капитализма к коммунизму, за которым позднее исторически закрепилось название «социализм», будет рассмотрена нами позднее. Пока речь идёт о том, к чему именно должен быть этот переход, или о том, что хотят получить коммунисты в результате революционных преобразований общественной жизни.

Нас здесь интересуют те обязательные определения, которые делают общество именно коммунистическим и отличают его от всех предшествующих коммунизму способов производства и осуществления общественной жизни. Поэтому мы оставим в стороне описание технологий при коммунизме. Они могут быть разными (и конечно же будет быстро изменяться). Например, капитализм при господстве парового двигателя и капитализм с современной цифровой техникой, очень сильно отличаются друг от друга. Но то, что делает капитализм капитализмом, обществом развитого товарного производства, где товаром является в том числе и рабочая сила, а капитал - господствующим общественным отношением - характерно, необходимо и неизменно для капитализма вообще, и в ХIX, и в ХXI веке. Понимание этих необходимых моментов в их необходимой взаимосвязи и составляет понятие капитализма.

Точно так же в понятие коммунизма, как и любого другого предмета, входит только то, что является необходимым и достаточным в своих необходимых взаимопереходах для того, чтобы общество было коммунистическим.

Трудность заключается в том, что мы имеем дело с предметом, ещё не ставшим реальностью. История ещё не знала (по крайней мере, на нашей планете) коммунистического общества, развитого на своём собственном основании (первобытный коммунизм в данном случае к делу не относится).

Однако, коммунизм уже давно стал наукой. Весь ход исторического развития мысли породил идею коммунизма, а ход общественной жизни, как минимум, постоянно показывает предел капитализма и необходимость выхода на новое основание    общественного     развития.    Капиталистические     кризисы перепроизводства сопровождаются колоссальным уничтожением не только произведенных товаров или отдельных предприятий, но часто и целых некогда процветающих мегаполисов.

Если рассматривать общественное производство как способ взаимодействия человека с природой, становится очевидно, что он нерационален и наносит нашей планете вред таких масштабов, что это может стать непоправимо. Производство и перепроизводство той же самой упаковки, излишней и бессмысленной с точки зрения разумного взаимодействия с природой, но необходимой для того, чтобы товар был продан - наносит огромный вред. Производство самих товаров, не обогащающих и не развивающих, а часто и калечащих человека, но поглощающих большое количество ресурсов, загрязняющее почву, воду и воздух - вот реальность капитализма. Отсутствие планомерной в масштабах всего общества, продуманной переработки отходов, а также разумного распределения и использования всех природных ресурсов и производственных возможностей. Всё это создаёт угрозу необратимой экологической катастрофы.

Попутно такому рукотворному разрушению окружающей нас природы, идёт деградация человека. Это не только абсолютная нищета, в которую ввергнута существенная часть трудящихся землян: люди голодают, живут в антисанитарных условиях, не имеют доступа к образованию и медицине; но речь также идёт об относительном обнищании, когда производятся новые потребности, удовлетворение которых становится необходимостью даже для того, чтобы воспроизводить себя в качестве продавца собственной рабочей силы, но средств для их удовлетворения недостаточно. Эту тенденцию отмечал ещё Энгельс, но в наши дни, когда капитализм не может жить без целенаправленного производства потребностей, она многократно усиливается и приобретает угрожающие масштабы, особенно на периферии и полупериферии капиталистического Мира, где в общественном масштабе рабочая сила не восстанавливается    в   полной    мере.   Это   подрывает     сами   основы капиталистического производства.

Но это ещё не всё. Современный капитал для того, чтобы воспроизводиться в расширенном масштабе возрождает и подчиняет себе докапиталистические способы эксплуатации человека человеком. Рабство и крепостничество в современном мире существуют постольку, поскольку они оказываются выгодными и способствуют самовозрастанию стоимости. Например, по данным МОТ, опубликованным в 2014 году доходы от использования рабского труда во всем мире превышают 150,2 миллиардов долларов (как видите, выгода от рабства исчисляется и может исчисляться чисто по-капиталистически, а не в натуральной форме, как это было в эпоху рабовладения).

Капитал постоянно требует простора для самовозрастания, а для этого ему нужно вытеснить другой капитал из уже имеющихся рынков, перераспределить ресурсы в свою пользу. Это порождает конфликты не только внутри государства, но, главным образом, между ними, что определяет локальные военные действия и ведёт к очередной мировой войне. А когда мы имеем тенденцию (подготовку, развёртывание, отдельные вспышки) мировой войны и ядерное оружие, какие только случайности не возможны!.. При этом сама война даёт мощный толчок развитию широкого, массового, стихийного антикапиталистического движения. Насколько последовательной будет эта линия, и будет ли она вести к реальному выходу за пределы капитализма - сказать пока трудно. Вообще гарантии, что человечество переживёт капитализм, никто сейчас дать не может. Единственное, о чём это свидетельствует, так это о пределе развития общества, где господствует капитал.

То, что капитализму придёт конец, несомненно. Вопрос только в том, закончится ли он вместе с человечеством или человечество научится жить по- другому. И как именно по-другому можно жить, чтобы и человечество в целом и все его многообразные составляющие развивались, не мешая, а помогая друг другу? Научная теория коммунизма даёт ответ на этот, отнюдь не устаревший, а животрепещущий для человечества вопрос.

Основные моменты этой теории мы и попытаемся здесь изложить. Она создавалась Марксом и Энгельсом. Внесли в развитие понимания коммунизма свою лепту и более поздние мыслители-марксисты. Среди них особого внимания заслуживает теоретическая работа Ленина.

Значение терминов «коммунизм» и «социализм» устоялось не сразу. Долгое время они употреблялись как синонимы или коммунизм даже могли понимать как предшествующий социализму. Потому следует наравне обращаться как к работам Маркса и Энгельса, где употребляется термин «коммунизм», так и к тем, где говорится о «социализме», если по сути изложенного ясно, что речь идёт о коммунистическом обществе, о коммунистическом движении или о теории коммунизма. Что же касается Ленина и других марксистов ХХ века, то в их работах, как и в поздних работах Маркса, терминология уже более-менее устоялась. Ленин в статье «Великий почин» специально напоминает читателям, что «научное различие между социализмом и коммунизмом только то, что первое слово означает первую ступень вырастающего из капитализма нового общества, второе слово — более высокую, дальнейшую ступень его». Там, где мы будем употреблять эти слова без ссылок и особых разъяснений, их следует понимать согласно современному словоупотреблению.

Насчет аргумента против коммунизма, на все лады повторяемого разного рода идейной обслугой современного правящего класса капиталистов, - «коммунизм - это утопия, уже пробовали построить коммунизм, и ничего (или ничего хорошего) из этого не получилось», - на это мы ответим следующее. Да, коммунистическое общество выстроить пока не получилось, но вот сказать, что не получилось ничего хорошего, - это попросту солгать. Успехи социалистического Советского Союза в свое время не могли отрицать даже самые рьяные антикоммунисты, даже в США. Те, кто помнит недавнюю историю ХХ века, хорошо знают, что направление движения общества к коммунизму позволило подняться на невиданные прежде, при феодализме и капитализме, высоты всем народам, жившим в СССР, причём в самые кратчайшие сроки. Об этом свидетельствуют как сухие цифры и выкладки, так и до сих пор не превзойденные достижения в области культуры.

Если сравнить первые десятилетия\footnote{Несмотря на то, что движение это начиналось в условиях разрушенной войной, по большей части аграрной страны, находящейся во враждебном окружении, было движением по ещё нехоженому пути и потому осуществлялось с большими проблемами, трудностями и отступлениями и даже трагедиями, о которых мы, коммунисты, не имеем права забывать. И происходило это движение в обществе, с довольно сильными капиталистическими тенденциями, которые нельзя отменить декретом или перешагнуть через них, не уничтожив товарное производство.} жизни общества, где было движение в сторону коммунизма и тот же самый срок движения «независимых государств» вразброд после распада СССР, мы обнаружим прямую зависимость между движением к коммунизму и развитием всего общества. Это касается не только и не столько промышленности и производства вещей, сколько возможностей для развития человека и человеческих способностей, доступных для широких масс населения.

Уже в 1925-м году в начальных школах училось на 2 250 000 учеников больше, чем в 1913-м. Число учащихся в профтехнических школах выросло вдвое. Из гражданской войны общество село сразу же за парту. Право получать образование было завоёвано для тех, кто до гражданской был обречён на невежество. Вдумайтесь: Советскому Союзу только год, а ему уже есть, что предъявить капиталистическому миру, в котором доминирует сословный принцип в образовании. Хорошее образование в капиталистическом мире было и остаётся недоступным для большинства.

За 20 лет после Октября удалось решить проблему беспризорности, причём лишившимся в ходе Первой Мировой и Гражданской родителей детям давался не просто кров и хлеб – давалось образование, профессии, путёвка в жизнь. В то время как на осколках бывшей единой социалистической страны снова появились беспризорники, причём по большей части родители их живы. Упал и продолжает падать уровень общей образованности населения. Системы образования бывших Советских республик стремятся к соответствию новым социальным реалиям: всё более узкой специализации и углубляющемуся расслоению в образовательных возможностях для детей бедных и богатых. Множество красноречивых свидетельств об этом процессе в России, вы найдёте в фильме Константина Сёмина и Евгения Спицина «Последний звонок».

В СССР программа ликвидации безграмотности (ликбез – с 1920-го) и всеобщего начального обучения (всеобуч – с 1930) смогла решить проблемы начального образования менее чем за двадцать лет – и это в аграрной стране, где до 1917-го года было более 80\% безграмотного и малограмотного населения. К 1936-му году в СССР было обучено около 40 млн неграмотных. В 1933—1937 годах только в учтённых школах ликбеза занимались свыше 20 млн неграмотных и около 20 млн малограмотных. По данным переписи 1939 года, грамотность лиц в возрасте от 16 до 50 лет приближалась к 90\%. Вот эти самые люди, а не мифические богатыри с крестиками на груди и богом в сердце (на чём настаивал фильм Никиты Михалкова «Цитадель» и другие изделия современной пропагандистской индустрии) и победили мощнейшую, самую современную и оснащённую на тот момент, военную машину нацистской Германии, на которую к моменту вторжения в СССР работала большая часть европейской индустрии.

Итак, за 20 лет была практически ликвидирована безграмотность, построена сеть мед. учреждений, где каждый имел доступ к медицинской помощи (врачи и учителя появились там, где их до этого никогда не было), появились школы, библиотеки, учреждения культуры, был дан мощный толчок для развития науки, в том числе и благодаря появлению соответствующей инфраструктуры. Строительство коммунизма как приоритетное всестороннее развитие человека (образование, здравоохранение, наука, культура) позволило победить угрожавший всей планете фашизм, а после тяжёлой разрушительной войны - в кратчайшие сроки отстроить страну.

Что же жизнь в годы после развала СССР принесла бывшим советским республикам полное сворачивание коммунистических тенденций в развитии общества?

В одной только Украине после распада СССР по данным Государственного комитета статистики до 2014 года абсолютная убыль населения составила почти 6,5 миллионов человек (до распада прирост был положительным). Этот показатель не учитывает тех, кто «числится» гражданами Украины, но вынужден был покинуть свой дом и поехать на заработки. Это само по себе свидетельствует о колоссальной социальной катастрофе. Конечно, тут вспомнят голод 1932-1933 годов. Массовый голод - всегда страшная трагедия, вне зависимости от её подлинных причин. Но даже по самым антисоветским подсчётам (именно подсчётам, а не безосновательным идеологическим заявлениям) проведенным Институтом демографии НАН Украины, голод на Украине унёс не 7 и не 10, а 3,9 миллиона человеческих жизней. Учитывались не просто потери, а сравнивалось количество населения с прогнозируемым количеством с учётом значительного положительного прироста в годы до и после голода. Если по аналогичной методике рассчитать абсолютные потери населения Украины после развала СССР, то картина будет гораздо страшнее, чем получившаяся в результате вычитания почти 6,5 млн.

В Литве, Латвии и Эстонии также фиксируется значительная депопуляция. В РФ за счёт миграции зафиксирована меньшая убыль населения с 148,6 на 1991 года до 144,3 млн. на 2016.

Если учесть трудовую миграцию в других направлениях и структуру занятости в самих среднеазиатских республиках, станет ясно, что общее качество жизни и возможность у граждан этих республик найти работу, позволяющую содержать семью, у себя дома стали значительно ниже.

Войны на территории бывшего СССР стали частым и обычным явлением: конфликт Азербайджана и Армении в Нагорном Карабахе (начавшийся ещё в СССР, но, по большому счёту, во время распада), гражданская война в Таджикистане, Приднестровье, две чеченские войны, Грузия и Абхазия, Южная Осетия, Донбасс. И во всех этих войнах люди убивали друг друга, в конечном счёте, в угоду борющимся группировкам капитала, то есть своим собственным эксплуататорам.

Нам могут возразить, что гражданская война, начавшаяся в 1917 году, тоже была кровопролитной и унесла, по различным подсчётам, от 8 до 13 млн жизней со всех сторон. Сюда относятся не только убитые в ходе боевых действий, но и умершие от болезней, голода и террора. Да, это так, ответим мы. И, конечно же, лучше было бы обойтись без войны. Но, во-первых, её развязали не большевики и не пролетариат, а представители старых эксплуататорских классов, не желающих лишаться своего господства. Во-вторых, в этой войне трудящиеся отстаивали свои собственные интересы, а не интересы тех, кто живёт эксплуатацией их труда. В-третьих, эта война была борьбой за возможность осуществить кардинальное преобразование всего общества, а не просто переделом сфер влияния капитала. О Великой Отечественной войне имеет смысл говорить отдельно, так как это было противостоянием фашизму, где Советский народ отстаивал не только свои интересы, но и интересы всего человечества.

Как видите, даже такое простое сравнение первых десятилетий существования СССР и жизни республик после его развала даёт возможность говорить о преимуществах общества, где есть простор для тенденции движения к коммунизму. И это при всём том, что процесс шёл совсем не гладко, и было много того, что уж никак нельзя назвать коммунистическим. Речь идёт о чисто капиталистических процессах и явлениях в социалистическом обществе. Именно о них так много кричат критики СССР, несмотря на то, что эта критика, там, где она вполне справедлива, оказывается самокритикой по отношению к капиталистическому миру.

Тенденция движения к коммунизму в СССР позволила всего через двенадцать лет после унесшей миллионы людских жизней и разрушившей существенную часть страны Великой Отечественной войны запустить первый искусственный спутник Земли. А вскоре и отправить первого человека в космос. Строительство коммунизма позволило разработать в научном плане те идеи, которые для нас в нашем нынешнем состоянии общества, скорее, являются идеями будущего, чем прошлого, так как их реализация не совместима с частной собственностью на средства производства. Список можно продолжать долго и включать в него не только советские достижения, но и достижения других социалистических стран, - возникших как раз на том пути, который виделся Ленину в свете гипотезы победы мировой революции.

Конечно же в этой грандиозной стройке были не только достижения, но и отступления, и провалы. Но, даже если не учитывать сказанного выше, аргумент

«не получилось построить коммунизм, значит, это вообще невозможно и ничего из этого не получится» страдает тем же логическим недостатком, каким страдали заявления противников идей, положенных, например, в основу появления и развития авиации. Противники утверждали, что невозможно построить летательные аппараты тяжелее воздуха. Это всё равно, что заявить: человек или человечество не сможет чему-то научиться только потому, что у него до сих пор это не получилось, хотя попытки были.

Такие нелогичные аргументы не предполагают изучения существа дела и возражений по существу. Поэтому и науки о коммунизме они совершенно не касаются. Неудача первого рывка человечества к коммунизму свидетельствует только о необходимости детального изучения исторического опыта для того, чтобы понять причину этой неудачи, более тщательного изучения условий, которые должны быть достигнуты для перехода к коммунизму. Само существование и состав коммунистического идеала, родившегося как осмысление тенденций и противоречий капиталистического общества, требует этого.

Итак, поскольку коммунизм, кроме всего прочего, стал наукой, приступим к его изучению. Брошюра эта, конечно же, не может претендовать на детальное и развернутое исследование коммунизма. Это, скорее, текст, предназначенный для того, чтобы изложить основные положения науки о коммунизме, разработанные в марксистской теории. Более глубокая и детальная проработка проблемы - дело ближайшего будущего. Мы приглашаем наших читателей присоединиться к этому делу, в том числе и в форме критики нашей работы.
\end{document}
